\section{计算与结果}
\label{sec:oneloopcalc}
\label{sec:oneloop}
虽然场的匹配系数可以是规范相关的，但需要注意，本工作所考虑的局域流算符所对应的组合是规范不变的。因此，我们在整个计算中始终采用费曼规范。我们的计算在欧氏时空中进行。在欧氏时空中，重正化的拉格朗日量 $\mathcal{L}_{h_v}$ 为
\begin{eqnarray}
\label{eq:extendLRenEuclidean}
\mathcal{L}_{h_v}^{\rm E} =Z_{h_v}\bar{h}_v (-iv\cdot \partial + i\delta m) h_v + gZ_gZ_A^{\frac{1}{2}}Z_{h_v} \bar{h}_v v\cdot A^{a}T^a h_v,
\end{eqnarray}
其中上标``E''表示欧氏时空。因此，
动量为 $k$ 沿``重''辅助场线流动的``重''辅助场传播子以及相互作用顶点的费曼规则分别为
$\frac{1}{k\cdot v -i0}$ 与 $-gv^{\mu}T^a$。欧氏时空中的 $-i0$ 规定对获得正确的线性发散和``质量修正''$i\delta m$ 起着关键作用。

我们按照如下方式把一个矢量分解为纵向与横向分量（在我们的设定下欧氏时空中 $v^2=1$）
\begin{eqnarray}
q_{\|}^{\mu} = (v\cdot q) v^{\mu}, \, \, \, \, q_{\perp}^{\mu} = q^{\mu} - v\cdot q v^{\mu}.
\end{eqnarray}
对 QCD 传播子分母，我们使用如下 Schwinger 参数化（$\alpha$ 表示）
\begin{eqnarray}
\label{eq:alphax}
\frac{1}{A^n} = \frac{1}{\Gamma(n)}\int_0^{+\infty} d\alpha\, \alpha^{n-1} e^{-A\alpha},
\end{eqnarray}
而对类晶（eikonal）传播子，我们使用另一版本的 Schwinger 参数化
\begin{eqnarray}
\label{eq:alphay}
\frac{1}{(k\cdot v- i0)^n} = \frac{(i)^n}{\Gamma(n)}\int_0^{+\infty} d\alpha\, \alpha^{n-1} e^{-i(k\cdot v)\alpha},
\end{eqnarray}
其中 $k\cdot v$ 取实值。


\subsection{$\delta m$ \& $\zeta_{h_v}$}
\begin{figure}
\centering
\includegraphics[width =0.32\textwidth]{images/Wilsonlineself.pdf}
\caption{梯度流形式下``重''辅助场的单圈自能费曼图。虚线箭头表示``重''辅助场，填充方块表示流时 $t$ 处流动的 $B_{\mu}$ 场。}
\label{fig:heavyself}
\end{figure}
梯度流中``重''辅助场单圈自能的计算与 HQET 中重夸克场的计算类似。
设外动量 $p$ 从左向右流动，图 \ref{fig:heavyself} 所示的单圈自能图给出振幅
\begin{eqnarray}
\mathcal{M}_{(a)} = g^2\tilde{\mu}^{2\epsilon}C_F\int\frac{d^dk}{(2\pi)^d}\frac{1}{\left(k\cdot v + p\cdot v - i0\right) k^2 } e^{-2tk^2},
\end{eqnarray}
其中 $\tilde{\mu}^{2\epsilon}= \left(\frac{\mu^2 e^{\gamma_E}}{4\pi} \right)^{\epsilon}$，且我们取``重''辅助场处于基础表示。对于伴随表示情形，只需在上式中把 $C_F$ 换成 $C_A$。

把``重''辅助场传播子的单圈自能图的几何级数求和得到 $1/(p\cdot v - \mathcal{M}_{(a)})$，这表明单圈``质量修正''由 $(-\mathcal{M}_{(a)})$ 给出。为简单起见，我们在 $\mathcal{M}_{(a)}$ 中取 $p=0$，得到
\begin{eqnarray}
\label{eq:deltamstart}
-\mathcal{M}_{(a)}|_{p=0}= -g^2\tilde{\mu}^{2\epsilon}C_F\int\frac{d^dk}{(2\pi)^d}\frac{1}{\left(k\cdot v- i0\right)k^2} e^{-2tk^2}.
\end{eqnarray}
注意上式在 $k\cdot v\rightarrow -k\cdot v$ 下并非奇函数，且 $k\cdot v$ 的积分并不给出为零的结果，因为 $i0$ 规定的符号至关重要。

为了在一圈推导``重''辅助场的匹配系数，我们对 $\mathcal{M}_{(a)}$ 关于 $p\cdot v$ 求导，并简单取 $p=0$，得到
\begin{eqnarray}
\label{eq:Zhvstart}
\frac{\partial\mathcal{M}_{(a)}}{\partial (p\cdot v)}|_{p=0}= - g^2\tilde{\mu}^{2\epsilon}C_F\int\frac{d^dk}{(2\pi)^d}\frac{1}{\left(k\cdot v- i0\right)^2k^2} e^{-2tk^2}.
\end{eqnarray}

现在，我们利用 eq.~(\ref{eq:alphax}) 与 eq.~(\ref{eq:alphay}) 中的 Schwinger 参数化显式求值 eq.~(\ref{eq:deltamstart}) 与 eq.~(\ref{eq:Zhvstart})。（当积分有限时我们取 $\epsilon=0$）我们发现
\begin{eqnarray}
\label{eq:deltamcalc0}
-\mathcal{M}_{(a)}|_{p=0} &=& -ig^2C_F\int\frac{d^4k}{(2\pi)^4}\int_0^{+\infty}dx\int_0^{+\infty}dy\, e^{-(2t+x)k^2- i(k\cdot v)y},\\
\label{eq:Zhvcalc0}
\frac{\partial\mathcal{M}_{(a)}}{\partial (p\cdot v)}|_{p=0} &=& g^2\tilde{\mu}^{2\epsilon}C_F\int\frac{d^dk}{(2\pi)^d}\int_0^{+\infty}dx\int_0^{+\infty}dy\, y\, e^{-(2t+x)k^2- i(k\cdot v)y},
\end{eqnarray}
完成 $k$ 动量积分后给出
\begin{eqnarray}
\label{eq:deltamcalc1}
-\mathcal{M}_{(a)}|_{p=0} &=& -i\frac{\alpha_s}{4\pi}C_F\int_0^{+\infty}dx\int_0^{+\infty}dy\, \frac{1}{\left(2t +x\right)^2}e^{-\frac{y^2}{4\left(2t+x\right)}},\\
\label{eq:Zhvcal1}
\frac{\partial\mathcal{M}_{(a)}}{\partial (p\cdot v)}|_{p=0} &=& \frac{g^2\tilde{\mu}^{2\epsilon}C_F}{(4\pi)^{d/2}} \int_0^{+\infty}dx\int_0^{+\infty}dy\, y\, \frac{1}{\left(2t +x\right)^{d/2}}e^{-\frac{y^2}{4\left(2t+x\right)}}.
\end{eqnarray}
积出 $x, y$ 后，我们得到结果
\begin{eqnarray}
\label{eq:deltamresult0}
-\mathcal{M}_{(a)}|_{p=0}&=& -i\frac{\alpha_s}{4\pi}C_F\frac{\sqrt{2\pi}}{\sqrt{t}},\\
\label{eq:Zhvresult0}
\frac{\partial\mathcal{M}_{(a)}}{\partial (p\cdot v)}|_{p=0} &=& - 2\times\frac{\alpha_sC_F}{4\pi}\left[\frac{1}{\epsilon_{\rm IR}}+\log\left(2\mu^2 te^{\gamma_E}\right)\right].
\end{eqnarray}

用取 $t=0$ 的维数正则化重复同样的计算，得到为零的 $\mathcal{M}_{(a),\, t=0}|_{p=0}$ 以及
\begin{eqnarray}
\frac{\partial\mathcal{M}_{(a),\, t=0}}{\partial (p\cdot v)}|_{p=0}=
2\frac{\alpha_sC_F}{4\pi}\left[\frac{1}{\epsilon_{\rm UV}}-\frac{1}{\epsilon_{\rm IR}}\right].
\end{eqnarray}
这证实了：当外动量取零时，从梯度流方案匹配到 $\overline{\rm MS}$ 方案只需从梯度流结果中减去 IR 极点。

因此，我们得到如下结果（注意重正化拉格朗日量中的``质量修正''记作 $i\delta m$）：
\begin{eqnarray}
\label{eq:deltamresult}
\delta m &=&  -\frac{\alpha_s}{4\pi}C_R\frac{\sqrt{2\pi}}{\sqrt{t}}+ O(\alpha_s^2),\\
\label{eq:Zhvresult}
\zeta_{h_v}  &=& 1- \frac{\alpha_sC_R}{4\pi}\log\left(2\mu^2 te^{\gamma_E}\right) + O(\alpha_s^2),
\end{eqnarray}
其中基础表示情形 $C_R=C_F$，伴随表示情形 $C_R=C_A$。上述 $\zeta_{h_v}$ 结果中对 $t$ 的对数依赖与文献 \cite{Craigie:1980qs} 中计算的 $h_v$ 重正化常数一致。

在下面的匹配计算中，我们略去未流动图的计算，因为从一开始就令 $t=0$ 便可以从相应流动图的振幅容易地得到它们，其结果正比于 $\frac{1}{\epsilon_{\rm UV}} - \frac{1}{\epsilon_{\rm IR}}$，显然与流动图的 IR 极点匹配。我们已经明确验证：下面及上面计算的匹配系数中 $-\log(2\mu^2te^{\gamma_E})$ 的系数与相应未流动算符单圈的 UV 极点系数一致，这可以作为我们计算的检验点。

%%%%%%%%%%%%%%%%%%%%%%%%%%%%%%%%%%%%%
\subsection{$\mathring{\zeta}_{\psi}$}
\begin{figure}
\centering
\includegraphics[width =0.72\textwidth]{images/quarkself.pdf}
\caption{梯度流形式下夸克单圈自能费曼图（镜像图隐含）。实线是费米子线，双实线是流费米子线，空心圆表示流顶点。}
\label{fig:quarkself}
\end{figure}
梯度流形式下夸克自能单圈修正的费曼图示于图 \ref{fig:quarkself}。类似的计算已在文献 \cite{Rizik:2020naq} 中完成；这里我们为完整性重做该计算，遵循把匹配系数拆分为若干可分别计算的部分的策略。

图 $(c_1)$ 在小流时极限下与普通 QCD 中的一样，其贡献为
\begin{eqnarray}
\label{eq:c1}
\frac{1}{i}\frac{\partial\mathcal{M}_{(c_1)}}{\partial \slashed{p}}|_{\slashed{p}=0}&=& -\frac{\alpha_sC_F}{4\pi}\left[\frac{1}{\epsilon_{\rm UV}}-\frac{1}{\epsilon_{\rm IR}}\right].
\end{eqnarray}
对图 $(c_2)$ 与 $(c_3)$，我们有
\begin{eqnarray}
\label{eq:c2}
\mathcal{M}_{(c_2)}&=& 2\times (i)^2g^2\tilde{\mu}^{2\epsilon}C_F\int_0^tds\, \int\frac{d^dk}{(2\pi)^d}\frac{(-2\slashed{k})\slashed{k}}{\left(k^2 \right)^2}e^{-sk^2}\nonumber\\
&=&2\times\frac{\alpha_sC_F}{4\pi}\left[\frac{1}{\epsilon_{\rm UV}}+1+\log\left(2\mu^2 te^{\gamma_E}\right)\right],
\end{eqnarray}
以及
\begin{eqnarray}
\label{eq:c3}
\mathcal{M}_{(c_3)}&=&2\times \frac{1}{2}g^2\tilde{\mu}^{2\epsilon}(-2dC_F)\int_0^tds\, \int\frac{d^dk}{(2\pi)^d}\frac{1}{k^2}e^{-2sk^2}\nonumber\\
&=&-4\times\frac{\alpha_sC_F}{4\pi}\left[\frac{1}{\epsilon_{\rm UV}}+\frac{1}{2}+\log\left(2\mu^2 te^{\gamma_E}\right)\right].
\end{eqnarray}

图 $(c_4)$ 及其对应的镜像图给出正比于外费米子动量的贡献（梯度流规范参数 $\kappa=1$），从而在外动量取零后消失。由于在胶子圈动量变换 $k\rightarrow -k$ 下固有的奇性，图 $(c_5)$ 给出零贡献。

把 $\frac{1}{i}\frac{\partial\mathcal{M}_{(c_1)}}{\partial \slashed{p}}|_{\slashed{p}=0}$、$\mathcal{M}_{(c_2)}$ 与 $\mathcal{M}_{(c_3)}$ 相加，通过引入带圈费米子场去除 UV 发散，再通过匹配手续减除 IR 发散，我们得到从带圈费米子场 $\mathring{\chi}$ 匹配到 $\overline{\rm MS}$ 重正化费米子场 $\psi$ 的单圈匹配系数：
\begin{eqnarray}
\label{eq:quarkselfresult}
\mathring{\zeta}_{\psi}= 1 +  \frac{1}{2} \times\frac{\alpha_s}{4\pi}C_F\left[\log\left(2\mu^2 te^{\gamma_E}\right)-\log(432)\right] + O(\alpha_s^2).
\end{eqnarray}
%%%%%%%%%%%%%%%%%%%%%%%%%%%%%%%%%%%%%%%%%%%%
\subsection{$\zeta_{\psi h_v}$}
\begin{figure}
\centering
\includegraphics[width =0.62\textwidth]{images/qvertex.pdf}
\caption{梯度流形式下 $\bar{\psi}h_v$ 顶点单圈修正的费曼图。叉圆表示 $\bar{\psi}h_v$ 顶点。连接叉圆的费米子场在流时 $t$ 流动。}
\label{fig:qvertex}
\end{figure}
梯度流形式下 $\bar{\psi}h_v$ 顶点的单圈修正费曼图由图 \ref{fig:qvertex} 给出。
把夸克外动量取为零，图 $(b_1)$ 的振幅为
\begin{eqnarray}
\label{eq:ZJstart}
\mathcal{M}_{(b_1)}= -g^2\tilde{\mu}^{2\epsilon}C_F\int\frac{d^dk}{(2\pi)^d}\frac{\left(\gamma\cdot v\right) \slashed{k}}{\left(-k\cdot v- i0\right)\left(k^2\right)^2} e^{-2tk^2}.
\end{eqnarray}
显然，$\slashed{k}$ 的横向分量给出为零的圈积分，于是我们有
\begin{eqnarray}
\label{eq:ZJresult0}
\mathcal{M}_{(b_1)}&=&g^2\tilde{\mu}^{2\epsilon}C_F\int\frac{d^dk}{(2\pi)^d}\frac{1}{\left(k^2\right)^2} e^{-2tk^2}\nonumber\\
&=&\frac{\alpha_sC_F}{4\pi}\left[-\frac{1}{\epsilon_{\rm IR}}-\log\left(2\mu^2 te^{\gamma_E}\right) -1\right].
\end{eqnarray}

图 $(b_2)$ 给出正比于外费米子动量的贡献（梯度流规范参数 $\kappa=1$），从而在外动量取零后消失。
因此，匹配到 $\overline{\rm MS}$ 方案后，我们得到 $\zeta_{\psi h_v}$ 的单圈结果：
\begin{eqnarray}
\label{eq:Zqhvresults}
\zeta_{\psi h_v}= 1 -\frac{\alpha_sC_F}{4\pi}\left[\log\left(2\mu^2 te^{\gamma_E}\right)+1\right] + O(\alpha_s^2).
\end{eqnarray}
%%%%%%%%%%%%%%%%%%%%%%%%%%%%%%%%%%%%%%%%%%%%%%%%
\subsection{$\zeta_A$}
\begin{figure}
\centering
\includegraphics[width =0.72\textwidth]{images/gluonself.pdf}
\caption{梯度流形式下胶子单圈自能的费曼图（镜像图隐含）。图 $(d_1)$ 中的圆团表示胶子真空极化。单曲线为胶子线，双曲线为流胶子线。}
\label{fig:gluonself}
\end{figure}
梯度流形式下胶子自能单圈修正的费曼图由图 \ref{fig:gluonself} 给出。为简单起见，提取相应贡献时我们把外动量取为零。图 $(d_1)$ 就是常规的胶子真空极化，它在小流时极限下给出的贡献为
\begin{eqnarray}
\label{eq:d1}
\mathcal{M}_{(d_1)}&=&\frac{\alpha_s}{4\pi}\delta^{\mu\nu}\left(\frac{5}{3}C_A - \frac{2}{3}n_f\right)\left[\frac{1}{\epsilon_{\rm UV}}-\frac{1}{\epsilon_{\rm IR}}\right].
\end{eqnarray}

图 $(d_2)$ 与 $(d_3)$ 给出的贡献为
\begin{eqnarray}
\label{eq:d2}
\mathcal{M}_{(d_2)}
&=& 4g^2C_A\tilde{\mu}^{2\epsilon}\int_0^t ds \int\frac{d^dk}{(2\pi)^d} \frac{(d-2)k^\mu k^\nu+k^2\delta^{\mu\nu}}{k^2 k^2}\text{exp}\left( -2s
k^2 \right)\nonumber\\
&=&\frac{\alpha_s}{4\pi}C_A\delta^{\mu\nu}\left[\frac{3}{\epsilon_{\rm UV}} + 3\log\left(2\mu^2 te^{\gamma_E}\right)+ \frac{5}{2}\right],
\end{eqnarray}
以及
\begin{eqnarray}
\label{eq:d3}
\mathcal{M}_{(d_3)}
&=& 2(1-d)g^2C_A\tilde{\mu}^{2\epsilon}\delta^{\mu\nu}\int_0^t ds \int\frac{d^dk}{(2\pi)^d} \frac{1}{k^2}\text{exp}\left( -2s
k^2 \right)\nonumber\\
&=&-\frac{\alpha_s}{4\pi}C_A\delta^{\mu\nu}\left[\frac{3}{\epsilon_{\rm UV}} + 3\log\left(2\mu^2 te^{\gamma_E}\right)+ 1\right],
\end{eqnarray}
由此我们看到图 $(d_2)$ 与 $(d_3)$ 中的 UV 发散相互抵消。

图 $(d_4)$ 给出的贡献为
\begin{eqnarray}
\label{eq:d4}
\mathcal{M}_{(d_4)}
&=& 4g^2C_A\tilde{\mu}^{2\epsilon}\int_0^t ds\int_0^s du \int\frac{d^dk}{(2\pi)^d}\frac{2(d-2)k^\mu k^\nu+\delta^{\mu\nu}k^2}{k^2 }e^{-2sk^2}\nonumber\\
&=& \frac{\alpha_s}{4\pi}C_A\delta^{\mu\nu}\left[\frac{2}{\epsilon_{\rm UV}} + 2 \log\left(2\mu^2 te^{\gamma_E}\right)  -\frac{1}{2} \right].
\end{eqnarray}

由于在小流时展开中给出正比于 $t$ 的贡献，图 $(d_5)$ 在小流时极限下不作贡献。$\mathcal{M}_{(d_2)}, \mathcal{M}_{(d_3)}$ 与 $\mathcal{M}_{(d_4)}$ 的结果与文献 \cite{Makino:2014taa} 中给出的一致。

把 $\mathcal{M}_{(d_1)}, \mathcal{M}_{(d_2)}, \mathcal{M}_{(d_3)}$ 与 $\mathcal{M}_{(d_4)}$ 相加，我们得到
\begin{eqnarray}
\label{eq:dsum}
\mathcal{M}_{(d)}
&=&\mathcal{M}_{(d_1)} + \mathcal{M}_{(d_2)} + \mathcal{M}_{(d_3)} + \mathcal{M}_{(d_4)} \nonumber\\
&=& \frac{\alpha_s}{4\pi}\delta^{\mu\nu}\left[\frac{1}{\epsilon_{\rm UV}} \beta_0 - \frac{1}{\epsilon_{\rm IR}} \left(\frac{5}{3}C_A - \frac{2}{3}n_f\right)+ 2 C_A\log\left(2\mu^2 te^{\gamma_E}\right)  +C_A\right],
\end{eqnarray}
其中 $\beta_0=\frac{11}{3}C_A - \frac{2}{3}n_f$。
IR 极点通过匹配手续去除，UV 极点则被吸收进 $\overline{\rm MS}$ 方案中强耦合的重正化
\begin{eqnarray}
Z_g = 1 - \frac{\alpha_s}{4\pi}\frac{1}{2\epsilon_{\rm UV}}\beta_0 + O(\alpha_s^2),
\end{eqnarray}
这与一般论断一致：$B_{\mu}$ 中的流胶子场不需要重正化，而 $B_{\mu}$ 中的强耦合需要重正化 \cite{Luscher:2011bx}。

最后，我们得到从流场 $B_{\mu}$ 匹配到 $\overline{\rm MS}$ 重正化胶子场 $gA_{\mu}$ 的匹配系数：
\begin{eqnarray}
\label{eq:ZAresults}
\zeta_A = 1+  \frac{\alpha_s}{4\pi}C_A\left[ \log\left(2\mu^2 te^{\gamma_E}\right)  +\frac{1}{2}\right] + O(\alpha_s^2).
\end{eqnarray}
注意，在单圈水平上 $\zeta_A$ 中的对数依赖与伴随表示下 $\zeta_{h_v}$（见 eq.~(\ref{eq:Zhvresult})）中的对数相消。然而，eq.~(\ref{eq:ZAresults}) 方括号中的常数项不被抵消。
%%%%%%%%%%%%%%%%%%%%%%%%%%%%%%%%%%%%%%%
\subsection{$\zeta_{\perp\perp}$ \& $\zeta_F$}
\begin{figure}
\centering
\includegraphics[width =0.8\textwidth]{images/gluonvertexcombine2.pdf}
\caption{在梯度流形式下对算符 $gF^{\mu\nu}_{\|\perp}h_v,\, gF^{\mu\nu}_{\perp\perp}h_v$ 与 $g\bar{h}_vF^{\mu\nu}_{\perp\perp}h_v$ 的单圈匹配有贡献的顶点型费曼图（镜像图隐含）。叉圆处与填充方块处的胶子场在流时 $t$ 流动。对于 $g\bar{h}_vF^{\mu\nu}_{\perp\perp}h_v$，为简洁起见图 $(e_1)-(e_5)$ 中未画出 $\bar{h}_v$ 场。$g\bar{h}_vF^{\mu\nu}_{\perp\perp}h_v$ 中的``重''辅助场处于基础表示，而 $gF^{\mu\nu}_{\|\perp}h_v,\, gF^{\mu\nu}_{\perp\perp}h_v$ 中的``重''辅助场处于伴随表示。图 $(e_6)$ 与 $(e_7)$ 分别只对算符 $g\bar{h}_vF^{\mu\nu}_{\perp\perp}h_v$ 与 $gF^{\mu\nu}_{\|\perp}h_v$ 给出贡献。}
\label{fig:gluonvertex}
\end{figure}
梯度流形式下算符 $gF^{\mu\nu}_{\|\perp}h_v,\, gF^{\mu\nu}_{\perp\perp}h_v$ 与 $g\bar{h}_vF^{\mu\nu}_{\perp\perp}h_v$ 的单圈顶点型费曼图示于图 \ref{fig:gluonvertex}。取外胶子动量为 $q$（流入顶点），我们有树图阶（LO）振幅
\begin{eqnarray}
\mathcal{M}_{F_{\|\perp}h_v}^{\rm LO} &=& ig\left(q_{\|}^{\mu}A^{\nu}_{\perp} -q_{_\perp}^{\nu}A^{\mu}_{\|} \right),\\
\mathcal{M}_{F_{\perp\perp }h_v}^{\rm LO} &=& ig\left(q_{\perp}^{\mu}A^{\nu}_{\perp} -q_{\perp}^{\nu}A^{\mu}_{\perp} \right),\\
\mathcal{M}_{\bar{h}_vF_{\perp\perp }h_v}^{\rm LO} &=& ig\left(q_{\perp}^{\mu}A^{\nu}_{\perp} -q_{\perp}^{\nu}A^{\mu}_{\perp} \right),
\end{eqnarray}
其中我们略去了色指标、丢掉了外``重''场，但出于匹配手续的方便保留了外胶子场。


由于我们只关心与外动量无关的匹配系数，对算符 $gF^{\mu\nu}_{\perp\perp}h_v$ 与 $g\bar{h}_vF^{\mu\nu}_{\perp\perp}h_v$ 的振幅可以取 $q\cdot v=0$，但对算符 $gF^{\mu\nu}_{\|\perp}h_v$ 的振幅则不行。算符 $gF^{\mu\nu}_{\perp\perp}h_v$ 与 $g\bar{h}_vF^{\mu\nu}_{\perp\perp}h_v$ 的单圈顶点修正的计算是直接而相似的。这三个算符共享相同的图 $(e_1) - (e_5)$。我们将只给出图 $(e_1) - (e_5)$ 对算符 $gF^{\mu\nu}_{\|\perp}h_v$ 与 $gF^{\mu\nu}_{\perp\perp}h_v$ 的振幅，并附上图 $(e_6)$ 对算符 $g\bar{h}_vF^{\mu\nu}_{\perp\perp}h_v$ 的振幅以及图 $(e_7)$ 对算符 $g\bar{h}_vF^{\mu\nu}_{\|\perp}h_v$ 的振幅。

对 $gF^{\mu\nu}_{\|\perp}h_v$ 的顶点修正计算并非那么直接，原因是它与规范非不变算符 $iv^{\mu}A_{\perp}^{\nu}(iv\cdot D - i\delta m) h_v$ 存在非平凡混合。
我们将在下一小节详细处理 $gF^{\mu\nu}_{\|\perp}h_v$。

图 $(e_1) - (e_6)$ 的振幅为
\begin{eqnarray}
\label{eq:e1}
\mathcal{M}_{(e_1)}
&=&-g^3C_A\tilde{\mu}^{2\epsilon}\int\frac{d^dk}{(2\pi)^d}\bigg\{ik^{\mu}A^{\alpha}\frac{\delta^{\nu\alpha}\left(-q-k\right)\cdot v+v^{\alpha}\left(2q-k\right)^{\nu}+v^{\nu}\left(2k-q\right)^{\alpha}}{(k\cdot v- i0)k^2(k-q)^2}\nonumber\\
&&\, \, \, \, \, \, \, \, \, \, \, \, \, \, \, \, \, \, \, \, \, \, \, \, \, \, \, \, \, \, \, \, \, \, \, \, \, \, \, \, \, \, \, \, \, \, \, \, \, \, \, \, \, \, \, \, \, \, \, \, \, \, \, \, \,  \, \, \, \, \, \, \, \, \, \, \, -\left(\mu\leftrightarrow\nu\right)\bigg\}e^{-t\left[k^2+(k-q)^2\right]},
\end{eqnarray}
\begin{eqnarray}
\label{eq:e2}
\mathcal{M}_{(e_2)}
&=&g^3C_A\tilde{\mu}^{2\epsilon}\int\frac{d^dk}{(2\pi)^d}\int_0^t ds\bigg\{ik^{\mu}\frac{v^{\alpha}\left(k-2q\right)^{\nu}+2\delta^{\alpha \nu}q\cdot v-2v^{\nu}(k-q)^{\alpha}}{(k\cdot v- i0)(k-q)^2}\nonumber\\
&&\, \, \, \, \, \, \, \, \, \, \, \, \, \, \, \, \, \, \, \, \, \, \, \, \, \, \, \, \, \, \, \, \, \, \, \, \, \, \, \, \, \, \, \, \, \, \, \, \, \, \, \, \, \, \, \, \, \, \, \, \, \, \, \, \, \, \, \, \,  \, \, \, \, \, \, \, \, \, \, \, \, \, \, \, \, -\left(\mu\leftrightarrow\nu\right)\bigg\}A^{\alpha}e^{-(s+t)(k-q)^2-(t-s)k^2},\\\nonumber\\
\label{eq:e3}
\mathcal{M}_{(e_3)}
&=&g^3C_A\tilde{\mu}^{2\epsilon}\int\frac{d^dk}{(2\pi)^d}\int_0^t ds\bigg\{ik^{\mu}\frac{\delta^{\nu\alpha}\left(q+k\right)\cdot v+2v^{\nu}(-k)^{\alpha}-2v^{\alpha}q^{\nu}}{(k\cdot v- i0)k^2}\nonumber\\
&&\, \, \, \, \, \, \, \, \, \, \, \, \, \, \, \, \, \, \, \, \, \, \, \, \, \, \, \, \, \, \, \, \, \, \, \, \, \, \, \, \, \, \, \, \, \, \, \, \, \, \, \, \, \, \, \, \, \, \, \, \, \, \, \, \, \, \, \, \, \, \, \, \, \, \, \, \, \,  -\left(\mu\leftrightarrow\nu\right)\bigg\}A^{\alpha}e^{-(t+s)k^2-(t-s)(k-q)^2},\\
\label{eq:e4}
\mathcal{M}_{(e_4)}
&=&-\frac{3}{2}g^3C_A\tilde{\mu}^{2\epsilon}\int\frac{d^dk}{(2\pi)^d } \frac{i\left(q^{\mu}A^{\nu}-q^{\nu}A^{\mu}\right)}{k^2(k-q)^2}e^{-t\left[k^2+(k-q)^2\right]},
\end{eqnarray}
\begin{eqnarray}
\label{eq:e5}
\mathcal{M}_{(e_5)}
&=&-g^3C_A\tilde{\mu}^{2\epsilon}\int\frac{d^dk}{(2\pi)^d }\int_0^t ds\frac{i(k+3q)^{\mu}A^{\nu}-i(k+3q)^{\nu}A^{\mu}}{k^2}e^{-(t+s)k^2-(t-s)(k-q)^2},\nonumber\\\\
\mathcal{M}_{(e_6)}
&=&g^3\left(C_F-\frac{C_A}{2}\right)\tilde{\mu}^{2\epsilon}\int\frac{d^dk}{(2\pi)^d }\int_0^t ds\frac{i\left(q^{\mu}_{\perp}A^{\nu}_{\perp}-q^{\nu}_{\perp}A^{\mu}_{\perp}\right)}{(k\cdot v-i0)^2k^2}e^{-2tk^2}.
\end{eqnarray}
提取 LO 表达式并在其后取 $q=0$，我们得到各图对算符 $gF^{\mu\nu}_{\perp\perp}h_v$ 的贡献：
\begin{eqnarray}
\label{eq:e1perpperpsimresult}
\zeta_{\perp\perp,\, (e_1)}&=&
-\frac{\alpha_sC_A}{4\pi}\times \frac{1}{2} \left[\frac{1}{\epsilon_{\rm IR}}+\log\left(2\mu^2 te^{\gamma_E}\right) +1\right],\\
\label{eq:e2perpperp}
\zeta_{\perp\perp,\, (e_2)}
&=&0,\\
\label{eq:e3perpperpresult}
\zeta_{\perp\perp,\, (e_3)}
&=&\frac{\alpha_s}{4\pi}C_A\times\frac{1}{16},\\
\label{eq:e4perpperpresult}
\zeta_{\perp\perp,\, (e_4)}&=&\frac{\alpha_sC_A}{4\pi}\times \frac{3}{2} \left[\frac{1}{\epsilon_{\rm IR}}+\log\left(2\mu^2 te^{\gamma_E}\right) +1\right],\\
\label{eq:e5perperpresult}
\zeta_{\perp\perp,\, (e_5)}
&=&- \frac{\alpha_sC_A}{4\pi}\times\frac{25}{16}.
\end{eqnarray}
把图 $(e_1) - (e_5)$ 的所有贡献相加并通过匹配手续丢弃 IR 极点，我们得到算符 $gF^{\mu\nu}_{\perp\perp}h_v$ 的顶点匹配系数：
\begin{eqnarray}
\label{eq:resultgluonvertexperpperp}
\zeta_{\perp\perp} &=& 1+  \frac{\alpha_sC_A}{4\pi}\left[\log\left(2\mu^2 te^{\gamma_E}\right) -\frac{1}{2}\right]+ O(\alpha_s^2).
\end{eqnarray}
类似地，我们有各图对算符 $g\bar{h}_vF^{\mu\nu}_{\perp\perp}h_v$ 的贡献：
\begin{eqnarray}
\zeta_{F,\, (e_1)}   &=& - \frac{\alpha_s}{4\pi}C_A\times \frac{1}{2}\left[\frac{1}{\epsilon_{\rm IR}} +\log\left(2\mu^2 te^{\gamma_E}\right) + 1\right],\\
\zeta_{F,\, (e_2)}   &=& 0,\\
\zeta_{F,\, (e_3)} &=&\frac{\alpha_s}{4\pi}C_A\times \frac{1}{16},\\
\zeta_{F,\, (e_4)}   &=&  \frac{\alpha_s}{4\pi}C_A\times \frac{3}{2}\left[\frac{1}{\epsilon_{\rm IR}} +\log\left(2\mu^2 te^{\gamma_E}\right) + 1\right],\\
\zeta_{F,\, (e_5)}   &=& - \frac{\alpha_s}{4\pi}C_A\times \frac{25}{16},\\
\zeta_{F,\, (e_6)}  &= & \frac{\alpha_s}{4\pi}\left(C_F -\frac{C_A }{2}\right)\times  2\left[\frac{1}{\epsilon_{\rm IR}} +\log\left(2\mu^2 te^{\gamma_E}\right)\right],
\end{eqnarray}
由此导出
\begin{eqnarray}
\label{eq:resultgluonvertexF}
\zeta_{\perp\perp}^F&= &1+ \frac{\alpha_s}{4\pi} \left[2C_F\log\left(2\mu^2 te^{\gamma_E}\right)+\frac{1}{2}C_A\right] + O(\alpha_s^2).
\end{eqnarray}
%%%%%%%%%%%%%%%%%%%%%%%%%%%%%%%%%%%%%%%%%%%%%%%%%%%%%%%
\subsection{$\zeta_{\|\perp}$}
图 $(e_4), (e_5)$ 没有混合，对 $gF^{\mu\nu}_{\perp\perp}h_v$ 与 $gF^{\mu\nu}_{\|\perp}h_v$ 给出相同的贡献：
\begin{eqnarray}
\label{eq:e4pperpresult}
\zeta_{\|\perp,\, (e_4)}&=&\frac{\alpha_sC_A}{4\pi}\times \frac{3}{2} \left[\frac{1}{\epsilon_{\rm IR}}+\log\left(2\mu^2 te^{\gamma_E}\right) +1\right],\\
\label{eq:e5pperpresult}
\zeta_{\|\perp,\, (e_5)}
&=&- \frac{\alpha_sC_A}{4\pi}\times\frac{25}{16}.
\end{eqnarray}
\\
复杂性来自图 $(e_1), (e_2), (e_3)$ 与 $(e_7)$，它们的振幅为
\begin{eqnarray}
\label{eq:e1pperp}
\mathcal{M}_{(e_1)}^{\|\perp}
&=&ig^3C_A\tilde{\mu}^{2\epsilon}\int\frac{d^dk}{(2\pi)^d}\left[\frac{(q+k)^{\mu}_{\|} A_{\perp}^{\nu}-(2q-k)^{\nu}_{\perp} A^{\mu}_{\|}}{k^2(k-q)^2} +\frac{v^{\mu}k_{\perp}^{\nu}(2k-q)_{\perp}\cdot A_{\perp} }{(k\cdot v- i0)k^2(k-q)^2} \right]\nonumber\\
&&\, \, \, \, \, \, \, \, \, \, \, \, \, \, \, \, \, \, \, \, \, \, \, \, \, \, \, \, \, \, \, \, \, \, \, \, \, \, \, \, \, \, \, \, \, \, \, \, \, \, \, \, \, \, \, \, \,  \, \, \, \, \, \, \, \, \, \, \, \,  \, \, \, \, \, \, \, \, \, \, \, \,  \, \, \, \, \times e^{-t\left[k^2+(k-q)^2\right]},
\end{eqnarray}
\begin{eqnarray}
\label{eq:e2pperp}
\mathcal{M}_{(e_2)}^{\|\perp}
&=&ig_s^3C_A\tilde{\mu}^{2\epsilon}\int\frac{d^dk}{(2\pi)^d}\int_0^t ds \left[\frac{2\left(q^{\mu}_{\|}A^{\nu}_{\perp} - q_{\perp}^{\nu}A_{\|}^{\mu}\right)}{(k-q)^2} +\frac{2v^{\mu}k_{\perp}^{\nu}(k-q)\cdot A}{(k\cdot v-i0)(k-q)^2}\right]\nonumber\\
&&\, \, \, \, \, \, \, \, \, \, \, \, \, \, \, \, \, \, \, \, \, \, \, \, \, \, \, \, \, \, \, \, \, \, \, \, \, \, \, \, \, \, \, \, \, \, \, \, \, \, \, \, \, \, \, \, \, \, \,  \, \, \, \, \, \, \, \, \, \, \,  \, \, \, \, \, \, \, \, \, \, \,  \, \, \, \, \times e^{-(s+t)(k-q)^2-(t-s)k^2},
\end{eqnarray}
\begin{eqnarray}
\label{eq:e3pperp}
\mathcal{M}_{(e_3)}^{\|\perp}
&=&ig_s^3C_A\tilde{\mu}^{2\epsilon}\int\frac{d^dk}{(2\pi)^d}\int_0^t ds \Bigg[\frac{(k+q)^{\mu}_{\|}A^{\nu}_{\perp}-(k+q)^{\nu}_{\perp}A^{\mu}_{\|}}{k^2}\nonumber\\
&&+\frac{(A\cdot v)\left(q^{\mu}_{\|}k^{\nu}_{\perp}-q^{\nu}_{\perp}k^{\mu}_{\|}\right)+2v^{\mu}k_{\perp}^{\nu}(k\cdot A)}{(k\cdot v-i0)k^2}\Bigg]e^{-(s+t)k^2-(t-s)(k-q)^2},
\end{eqnarray}
\begin{eqnarray}
\label{eq:e7pperp}
\mathcal{M}_{(e_7)}^{\|\perp}
&=&ig^3C_A\tilde{\mu}^{2\epsilon}\int\frac{d^dk}{(2\pi)^d}\frac{-v^{\mu}A_{\perp}^{\nu} }{(k\cdot v-i0)(k-q)^2}e^{-2t(k-q)^2}.
\end{eqnarray}
即使在对圈积分做张量约化之后，这些振幅也不给出严格正比于算符 LO 表达式 $gF^{\mu\nu}_{\|\perp}h_v$ 的表达式。它们导致与规范非不变算符 $iv^{\mu}A_{\perp}^{\nu}(iv\cdot D - i\delta m) h_v$ 的混合，后者的 LO 表达式正比于 $q_{\|}^{\mu}A_{\perp}^{\nu}$，因此不容易把这种混合与那些对算符 $gF^{\mu\nu}_{\|\perp}h_v$ 给出贡献的部分区分开。幸运的是，我们可以通过与 $q_{\perp}^{\nu}A_{\|}^{\mu}$ 成正比的结果提取对算符 $gF^{\mu\nu}_{\|\perp}h_v$ 的贡献。为检验文献 \cite{Dorn:1981wa,Zhang:2018diq,Wang:2019tgg} 中发现的混合模式与线性发散的抵消，我们在小流时极限下也保留与 $q_{\|}^{\mu}A_{\perp}^{\nu}$ 和 $v_{\|}^{\mu}A_{\perp}^{\nu}/\sqrt{t}$ 成正比的结果。

如果在 $\mathcal{M}_{(e_1)}^{\|\perp}$ 与 $\mathcal{M}_{(e_7)}^{\|\perp}$ 中取 $t=0$，容易看到来自 $\mathcal{M}_{(e_7)}^{\|\perp}$ 与 $\mathcal{M}_{(e_1)}^{\|\perp}$ 第二项的线性发散在 $d=3-2\epsilon$ 时相互抵消。然而在梯度流方案中，线性发散的抵消要复杂得多。显然，当 $d=4-2\epsilon$ 且 $t>0$ 时，来自 $\mathcal{M}_{(e_7)}^{\|\perp}$ 与 $\mathcal{M}_{(e_1)}^{\|\perp}$ 第二项的线性发散并不抵消。只有把图 $(e_2)$ 与 $(e_3)$ 包括进来时，梯度流方案中线性发散的抵消才得以恢复。

通过采用 HQET 匹配计算中成熟的方法，可以为我们的匹配计算带来极大的简化。也就是说，我们聚焦大圈动量区域 $k >> q$，并在做圈动量积分之前把振幅展开到外动量 $q$ 的线性阶。
该方法在我们情形下有效，因为在小流时极限下流动算符与相应的未流动算符具有相同的 IR 行为，且对匹配系数的贡献来自 UV 区域。

按 $k >> q$ 把振幅展开至 $q$ 的线性阶并完成圈积分，我们得到
\begin{eqnarray}
\label{eq:e1pperpexpand}
\mathcal{M}_{(e_1)}^{\|\perp}
&=& \frac{\alpha_s}{4\pi}gC_A\left\{ -\frac{2}{3}\frac{\sqrt{2\pi}}{\sqrt{t}}v^{\mu}A_{\perp}^{\nu} - i \ell_1 \left[\left(\frac{3}{2} + \frac{4}{d}\right)q^{\mu}_{\|} A_{\perp}^{\nu}-\frac{3}{2}q^{\nu}_{\perp} A^{\mu}_{\|}\right]\right\},
\end{eqnarray}
\begin{eqnarray}
\label{eq:e2pperpexpand}
\mathcal{M}_{(e_2)}^{\|\perp}
&=& \frac{\alpha_s}{4\pi}gC_A\left\{ -\frac{1}{6}\frac{\sqrt{2\pi}}{\sqrt{t}}v^{\mu}A_{\perp}^{\nu} + i \left[\frac{15}{8}q^{\mu}_{\|} A_{\perp}^{\nu}-\frac{1}{8}q^{\nu}_{\perp} A^{\mu}_{\|}\right]\right\},
\end{eqnarray}
\begin{eqnarray}
\label{eq:e3pperpexpand}
\mathcal{M}_{(e_3)}^{\|\perp}
&=& \frac{\alpha_s}{4\pi}gC_A\left\{ -\frac{1}{6}\frac{\sqrt{2\pi}}{\sqrt{t}}v^{\mu}A_{\perp}^{\nu} + i \left[\frac{11}{16}q^{\mu}_{\|} A_{\perp}^{\nu}-\frac{15}{16}q^{\nu}_{\perp} A^{\mu}_{\|}\right]\right\},
\end{eqnarray}
\begin{eqnarray}
\label{eq:e7pperpexpand}
\mathcal{M}_{(e_7)}^{\|\perp}
&=& \frac{\alpha_s}{4\pi}gC_A\left\{\frac{\sqrt{2\pi}}{\sqrt{t}}v^{\mu}A_{\perp}^{\nu} + \frac{i}{2} \ell_1\frac{1}{2}q^{\mu}_{\|} A_{\perp}^{\nu}\right\},
\end{eqnarray}
其中
\begin{eqnarray}
\ell_1 = \frac{1}{\epsilon_{\rm IR}}+\log\left(2\mu^2 te^{\gamma_E}\right) +1.
\end{eqnarray}
从上述结果可以清楚地看到线性发散的抵消。减去对算符 $gF^{\mu\nu}_{\|\perp}h_v$ 的贡献之后，遗留的对 $t$ 的对数依赖也与文献 \cite{Dorn:1981wa,Zhang:2018diq,Wang:2019tgg} 中给出的混合的 UV 发散一致。

总结图 $(e_1) - (e_5)$ 与 $(e_7)$ 的所有贡献，我们得到算符 $gF^{\mu\nu}_{\|\perp}h_v$ 顶点修正贡献的匹配系数
\begin{eqnarray}
\label{eq:resultgluonvertexpperp}
\zeta_{\|\perp} &=& 1 -\frac{1}{2}\frac{\alpha_s}{4\pi}C_A+ O(\alpha_s^2),
\end{eqnarray}
它抵消了 eq.~(\ref{eq:ZAresults}) 中 $\zeta_{A}$ 方括号里的常数项，从而使对算符 $gF^{\mu\nu}_{\|\perp}h_v$ 匹配系数的单圈贡献为零。\footnote{在文献 \cite{Eller:2021qpp} 中，作者在梯度流形式下单圈计算了色电场关联函数 $G_E$。小流时极限下 $G_E$ 的匹配计算可以约化为对局域流算符 $g\bar{h}_vF^{\mu\nu}_{\|\perp}h_v$（$h_v$ 处于基础表示）的匹配计算，后者对匹配系数的单圈贡献同样为零。}


%%%%%%%%%%%%%%%%%%%%%%%%%%%%%%%%%%%%%%%%%%%%%%%%%%%%%%%
\subsection{威尔逊线算符的结果与应用}
我们现在准备总结四个局域流算符 $\bar{\psi} h_v$、$gF^{\mu\nu}_{\|\perp}h_v$、$gF^{\mu\nu}_{\perp\perp}h_v$ 与 $g\bar{h}_vF^{\mu\nu}_{\perp\perp}h_v$ 匹配系数的最终结果。局域流算符匹配系数与场、顶点匹配系数之间的关系见 eqs.~(\ref{eq:operatorrelations})。
$\zeta_{h_v}$、$\mathring{\zeta}_{\psi}$、$\zeta_{\psi h_v}$、$\zeta_{A}$、$\zeta_{\|\perp}$、$\zeta_{\perp\perp}$ 与 $\zeta_F$ 的单圈结果分别在 eq.~(\ref{eq:Zhvresult})、eq.~(\ref{eq:quarkselfresult})、eq.~(\ref{eq:Zqhvresults})、eq.~(\ref{eq:ZAresults})、eq.~(\ref{eq:resultgluonvertexpperp})、eq.~(\ref{eq:resultgluonvertexperpperp}) 与 eq.~(\ref{eq:resultgluonvertexF}) 中给出。代入这些结果，最后我们得到
\begin{eqnarray}
\label{eq:cpsihvfinal}
c_{\psi h_v}(t, \mu) &=& 1- \frac{\alpha_s}{4\pi}C_F\left[\frac{3}{2}\log\left(2\mu^2 te^{\gamma_E}\right)+\frac{\log(432)}{2}+1\right] +O(\alpha_s^2),\\
\label{eq:cpperpfinal}
c_{\|\perp} (t, \mu) &=& 1+O(\alpha_s^2),\\
\label{eq:cperpperpfinal}
c_{\perp\perp} (t, \mu) &=&1+ \frac{\alpha_s}{4\pi}C_A\times\log\left(2\mu^2 te^{\gamma_E}\right) +O(\alpha_s^2),\\
\label{eq:cpperpFfinal}
c_F(t, \mu)  &=& 1+\frac{\alpha_s}{4\pi}C_A\times\log\left(2\mu^2 te^{\gamma_E}\right) +O(\alpha_s^2).
\end{eqnarray}
我们已验证 eq.~(\ref{eq:cpsihvfinal})、eq.~(\ref{eq:cpperpfinal}) 与 eq.~(\ref{eq:cperpperpfinal}) 中对 $t$ 的对数依赖分别与文献 \cite{Dorn:1980hs,Dorn:1981wa,Dorn:1986dt} 给出的流算符 $\bar{\psi} h_v$、$gF^{\mu\nu}_{\|\perp}h_v$ 与 $gF^{\mu\nu}_{\perp\perp}h_v$ 的反常量纲一致（比较可参见采用类似记号的文献 \cite{Braun:2020ymy}）。对流算符 $g\bar{h}_vF^{\mu\nu}_{\perp\perp}h_v$，其反常量纲也在文献 \cite{Eichten:1990vp,Falk:1990pz} 中用背景场方法计算过（背景场方法的综述见文献 \cite{Abbott:1981ke}），与 eq.~(\ref{eq:cpperpFfinal}) 中的对数依赖一致。值得注意的是，虽然 $gA$ 在背景场方法中不对 $g\bar{h}_vF^{\mu\nu}_{\perp\perp}h_v$ 的反常量纲作贡献，但在本文采用的方法中它却有贡献。图 $(e_1)$ 与 $(e_4)$ 在 $\overline{\rm MS}$ 方案中的 UV 极点依计算是否采用背景场方法而不同，原因在于三胶子顶点的费曼规则不同。我们已验证这些差异恰好补偿了采用或不采用背景场方法（费曼规范）时 $gA$ 反常量纲带来的差异。有了局域流算符匹配系数以及``重''辅助场 $h_v$、夸克场 $\psi$ 和 $gA$ 对 $t$ 的正确对数依赖，相应顶点匹配系数的对 $t$ 对数依赖便得到了保证。虽然 $c_{\perp\perp}(t, \mu)$ 与 $c_F(t, \mu)$ 在单圈水平上一致，有趣的是在两圈水平上 $gF^{\mu\nu}_{\perp\perp}h_v$ 与 $g\bar{h}_vF^{\mu\nu}_{\perp\perp}h_v$ 的反常量纲相差一个正比于 $\pi^2$ 的项（两圈反常量纲见文献 \cite{Braun:2020ymy, Amoros:1997rx,Czarnecki:1997dz}）。此外，对算符 $gF^{\mu\nu}_{\|\perp}h_v$，其反常量纲从两圈开始出现 \cite{Braun:2020ymy}；算符 $\bar{\psi} h_v$ 的两圈与三圈反常量纲分别见文献 \cite{Broadhurst:1991fz,Ji:1991pr} 与文献 \cite{Chetyrkin:2003vi}。

如第 \ref{subsec:smallflowexp} 节末所述，在一维辅助场形式中，小流时极限下威尔逊线算符的匹配系数可以直接由相应局域流算符的匹配系数得到。在小流时极限下，把 eq.~(\ref{eq:defqoperator})、eq.~(\ref{eq:defpperp}) 与 eq.~(\ref{eq:defperpperp}) 定义的威尔逊线算符的匹配方程写为
\begin{eqnarray}
\mathcal{O}_{\Gamma}^{\rm R}(zv, t) &=& \mathcal{C}_{\psi}(t, \mu)e^{\delta m z}\mathcal{O}_{\Gamma}^{\overline{\rm MS}}(zv) + O(t) ,\\
\mathcal{O}_{\|\perp}^{\rm R} (zv, t) &=&  \mathcal{C}_{\|\perp}(t, \mu)e^{\delta m z}\mathcal{O}_{\|\perp}^{\overline{\rm MS}} (zv)+ O(t),\\
\mathcal{O}_{\|\perp}^{\rm R} (zv, t) &=& \mathcal{C}_{\perp\perp}(t, \mu)e^{\delta m z}\mathcal{O}_{\perp\perp}^{\overline{\rm MS}} (zv)+ O(t),
\end{eqnarray}
其中为简单起见我们省略了洛伦兹指标，由此我们建立关系
\begin{eqnarray}
\mathcal{C}_{\psi}(t, \mu) &=& c^2_{\psi h_v}(t, \mu),\\
\mathcal{C}_{\|\perp}(t, \mu) &=& c^2_{\|\perp} (t, \mu),\\
\mathcal{C}_{\perp\perp}(t, \mu) &=& c^2_{\perp\perp} (t, \mu).
\end{eqnarray}
这些关系与 $z$ 以及用这些威尔逊线算符构造的矩阵元的外态无关。

关于以色磁场插入威尔逊圈定义的自旋依赖势，每一次这样的插入都会在从梯度流方案到 $\overline{\rm MS}$ 方案的匹配中带来一个匹配系数 $c_F(t, \mu)$ 因子。
